\subsection{Subgroup Vulnerability and Health Inequity}

The impact of behavioral delay is not shared equally across the population. Our simulations show that subgroups with higher $\kappa$ values experience significantly more severe outbreaks. The most pronounced effect is observed in the \textit{Neutral Institutional Trust} subgroup ($\kappa = 0.356$), where the final Attack Rate reaches 65.88\%. 

Crucially, our findings highlight a "double burden" of risk for vulnerable populations. Individuals with \textit{chronic medical conditions} are not only at higher biological risk for severe outcomes but also face an elevated behavioral hurdle ($\kappa = 0.288$), leading to a cumulative Attack Rate of 64.70\%. Similarly, those with \textit{prior adverse vaccine reactions} ($\kappa = 0.315$) and \textit{adults aged 65+} ($\kappa = 0.327$) exhibit substantial resistance to administrative delay, resulting in Attack Rates that exceed 65.1\%. 

These results reveal that administrative friction functions as a regressive tax on public health. Those who are most biologically vulnerable or most skeptical of institutional delivery are precisely the ones most likely to be deterred by service delays. Consequently, a "one-size-fits-all" vaccination rollout that fails to streamline access for these high-$\kappa$ groups will inadvertently exacerbate health disparities, concentrating the burden of the epidemic among the most marginalized and at-risk segments of society.
